%% sections/slide06-other-results.tex — SLIDE 6 — All other results
%% =============================================================================
%% CONTENT BLUEPRINT — Secondary results, compact
%% =============================================================================
%% Time budget: 60 seconds.
%% Cram as many supporting findings as fit. Use a 2×2 figure grid with
%% one-line captions per panel. This is the data wall.
%%
%% WHAT TO PUT HERE:
%%
%%   Layout: 2 rows × 2 columns mini-figure grid. Each panel gets:
%%     • A small plot (4–5 cm tall, full-width within its cell).
%%     • A header line stating the result: "R2: E_g decreases 8 % per
%%       1 % tensile strain."
%%     • No prose inside the panel — keep panels tight.
%%
%%   Below the grid (optional): a 1-line cross-reference to data archived
%%   at <DOI / GitHub> for the examiner to dig in if interested.
%%
%% FIGURE SELECTION GUIDELINES:
%%   • Each figure must answer a specific question from slide 3.
%%   • If a finding is too small for its own panel, drop it. Better to
%%     have three strong figures than four crammed ones.
%%   • Order: figures must approximately map to the order of objectives
%%     O2 → O3 → O4 from slide 3, so the examiner can trace.
%%
%% ANNOTATIONS:
%%   • Same rule as slide 5: annotate the most relevant feature inline
%%     inside each panel (arrow + label).
%%   • No axis labels saying "x" / "y" — use the actual quantity.
%%
%% IF YOU HAVE >4 RESULTS:
%%   • Push the least essential to the thesis (don't put everything in
%%     slides — the examiner will read the thesis afterwards).
%%   • Or split into slides 6a + 6b, but timer considers this a sign you
%%     have too weak a headline. Strengthen slide 5 instead.
%% =============================================================================

\begin{frame}{Other Results~\textemdash~supporting findings}
  \begin{itemize}\setlength\itemsep{1pt}
    \item \highlight{\textbf{R2}} \, [Second-most-important finding, single
          sentence with number and conditions.]
    \item \highlight{\textbf{R3}} \, [Third finding, single sentence with
          number and conditions.]
    \item \highlight{\textbf{R4}} \, [Optional: fourth finding, single
          sentence with number and conditions.]
  \end{itemize}

  \vspace{0.15cm}
  \begin{columns}[T, totalwidth=\textwidth]
    \begin{column}{0.48\textwidth}
      \centering
      \vspace{-2pt}
      \textbf{\color{sustBlueDark}R2 — growth / fitting / trend}\\[-2pt]
      \scriptsize Single-line take-away sentence ends with a number.
      \begin{tikzpicture}
        \begin{axis}[
            width=\textwidth, height=3.6cm,
            xlabel={\tiny $x$}, ylabel={\tiny $Y$},
            xmin=0, xmax=10,
            grid=major, grid style={line width=.1pt, color=gray!20},
            tick label style={font=\tiny}
          ]
          \addplot[sustBlue, very thick, mark=square*, mark size=1pt] coordinates {
            (1,1.0)(2,1.4)(3,1.7)(4,1.9)(5,2.0)(6,1.9)(7,1.7)(8,1.4)(9,1.0)};
        \end{axis}
      \end{tikzpicture}
    \end{column}

    \begin{column}{0.48\textwidth}
      \centering
      \vspace{-2pt}
      \textbf{\color{sustBlueDark}R3 — critical comparison}\\[-2pt]
      \scriptsize Single-line take-away sentence ends with a number.
      \begin{tikzpicture}
        \begin{axis}[
            width=\textwidth, height=3.6cm,
            xlabel={\tiny $x$}, ylabel={\tiny $Y$},
            xmin=0, xmax=10,
            grid=major, grid style={line width=.1pt, color=gray!20},
            tick label style={font=\tiny}
          ]
          \addplot[sustAccent, very thick, mark=*, mark size=1pt] coordinates {
            (1,2.1)(2,1.7)(3,1.4)(4,1.2)(5,1.1)(6,1.2)(7,1.3)(8,1.5)(9,1.8)};
        \end{axis}
      \end{tikzpicture}
    \end{column}
  \end{columns}

  \vspace{0.1cm}
  \timenote{Full data + scripts deposited: \texttt{[DOI / GitHub link]}.}
\end{frame}
